\section{Estudio de alternativas}
\label{sec:alternativas}

Este capítulo analiza las principales alternativas tecnológicas consideradas para cada componente del sistema y justifica las elecciones adoptadas. La selección de herramientas en un proyecto de este tipo no es trivial: los ecosistemas de orquestación de agentes, proveedores de modelos y almacenes de datos han experimentado una expansión acelerada en los últimos años, y no todas las opciones ofrecen el mismo grado de control, reproducibilidad o adecuación para un prototipo de investigación académica. En cada subsección se identifican los criterios relevantes para la decisión, se comparan las alternativas y se razona la elección final.


\subsection{Alternativas a LangGraph como framework de orquestación}
\label{sec:alt-langgraph}

La elección del framework de orquestación es la decisión arquitectónica más determinante del proyecto: condiciona la forma en que se modelan los estados, los nodos, el flujo de control y el mecanismo de supervisión humana. Se evaluaron cuatro familias de alternativas.

\textbf{AutoGen} (Microsoft) organiza los agentes como participantes en conversaciones de múltiples turnos~\cite{wu2023autogen}. Es adecuado para flujos conversacionales abiertos pero resulta inadecuado para pipelines con estado tipado, transiciones condicionales deterministas y puntos de interrupción con persistencia. AutoGen no distingue nodos deterministas de nodos agénticos ni proporciona un mecanismo nativo de checkpointing que permita suspender y reanudar la ejecución.

\textbf{CrewAI} ofrece una abstracción de alto nivel basada en roles (\emph{Crew}, \emph{Agent}, \emph{Task}). Permite construir sistemas multi-agente con poca configuración, lo que lo hace atractivo para prototipos simples. Sin embargo, el nivel de abstracción dificulta el control granular del flujo: no es posible definir nodos puramente deterministas, mezclar código Python con agentes LLM en la misma infraestructura, ni implementar un patrón HITL mediante interrupción del grafo.

\textbf{Semantic Kernel} (Microsoft) es un SDK orientado a la integración de LLMs mediante \emph{plugins} y funciones, con buen soporte para escenarios empresariales en C\# y Python. Su abstracción principal no es el grafo de estados sino el pipeline de funciones encadenadas, lo que dificulta la implementación de ciclos, bifurcaciones condicionales y estados compartidos con semántica de reducer.

\textbf{Apache Airflow} y herramientas equivalentes de orquestación de tareas (Prefect, Dagster) gestionan DAGs de tareas con dependencias, persistencia y monitorización. Son soluciones maduras para pipelines ETL y MLOps deterministas, pero no están diseñadas para integrar razonamiento LLM en nodos del grafo ni para implementar HITL mediante suspensión del estado en tiempo de ejecución.

\textbf{Código Python sin framework} ofrece máximo control pero requiere implementar desde cero el estado compartido, los reducers de acumulación, el checkpointer de persistencia y el mecanismo de interrupción para HITL. El coste de desarrollo de esta infraestructura supera el valor que aporta en el contexto de un TFG cuyo objetivo principal no es el framework de orquestación.

\begin{figure}[H]
\centering
\begin{tabular}{lcccc}
\hline
\textbf{Criterio} & \textbf{LangGraph} & \textbf{AutoGen} & \textbf{CrewAI} & \textbf{Airflow} \\
\hline
Grafos con ciclos y estado tipado & Sí & No & Parcial & No \\
Nodos deterministas y agénticos unificados & Sí & No & No & Parcial \\
Checkpointer nativo (HITL) & Sí & No & No & No \\
Control granular del flujo condicional & Sí & Parcial & No & Sí \\
Integración nativa con LangChain & Sí & No & Parcial & No \\
\hline
\end{tabular}
\caption{Comparativa de frameworks de orquestación de agentes.}
\label{tab:alt-langgraph}
\end{figure}

\textbf{Justificación de la elección.} LangGraph es el único framework evaluado que satisface simultáneamente los tres requisitos estructurales del proyecto: (1) estado compartido tipado con semántica de actualización parcial, (2) tratamiento uniforme de nodos deterministas y agénticos, y (3) mecanismo nativo de interrupción y reanudación del grafo para HITL. Ninguna de las alternativas ofrece esta combinación sin requerir desarrollo de infraestructura adicional.


\subsection{Alternativas al proveedor del modelo de lenguaje}
\label{sec:alt-proveedor}

Los nodos agénticos del sistema requieren un modelo de lenguaje capaz de seguir instrucciones complejas, producir salidas estructuradas validadas por Pydantic y mantener coherencia a lo largo de contextos de varios miles de tokens. La elección del proveedor afecta al coste, la latencia, la fiabilidad y la capacidad de seguimiento de instrucciones.

\textbf{API de OpenAI (acceso directo)} ofrece los modelos \texttt{gpt-4.1-mini} y \texttt{gpt-4.1} con alta fiabilidad, excelente capacidad de seguimiento de instrucciones y sin restricciones de tasa. Su principal limitación para un proyecto académico es el coste: a razón de entre 0,15 y 15 USD por millón de tokens según el modelo, el coste de un ciclo completo de desarrollo y pruebas puede ser significativo.

\textbf{GitHub Models} proporciona acceso gratuito a los mismos modelos de OpenAI (entre otros) a través de un punto de acceso compatible con la API de OpenAI~\cite{github_models}. La autenticación se realiza mediante un token personal de GitHub (\texttt{GITHUB\_TOKEN}), lo que simplifica la configuración sin requerir cuenta de pago. La restricción principal es una cuota de 150 peticiones por día en el nivel gratuito, lo que limita la velocidad de desarrollo iterativo pero es suficiente para la ejecución del pipeline en sesiones de trabajo individuales.

\textbf{Anthropic (Claude)} ofrece modelos con capacidades competitivas para salida estructurada y razonamiento encadenado. Requiere el SDK de Anthropic, cuya API difiere de la de OpenAI, lo que implicaría modificar la factoría de modelos y las integraciones de LangChain. La migración es técnicamente factible pero añade complejidad sin beneficio diferencial para este caso de uso.

\textbf{Ollama (modelos locales)} permite ejecutar modelos como Llama 3, Mistral o Phi-3 en local sin coste por petición ni dependencia de internet. Es la opción preferida cuando la privacidad de los datos es un requisito o cuando se necesita eliminar la dependencia de un proveedor externo. Sin embargo, los modelos locales de tamaño moderado (7B--13B parámetros) presentan mayor variabilidad en el seguimiento de instrucciones estructuradas y en la producción de JSON válido, lo que introduce inestabilidad en nodos como el planificador, donde la validación de cuatro etapas del contrato Pydantic requiere alta fidelidad.

\textbf{Hugging Face Inference API} ofrece acceso a cientos de modelos abiertos mediante una API REST. La variabilidad en la calidad de seguimiento de instrucciones entre modelos dificulta la selección y el mantenimiento de un conjunto estable de nodos agénticos.

\begin{figure}[H]
\centering
\begin{tabular}{lccccc}
\hline
\textbf{Criterio} & \textbf{GitHub Models} & \textbf{OpenAI directo} & \textbf{Anthropic} & \textbf{Ollama} \\
\hline
Compatible con API OpenAI & Sí & Sí & No & Parcial \\
Coste cero para prototipado & Sí & No & No & Sí \\
Capacidad structured output & Alta & Alta & Alta & Media \\
Sin dependencia de internet & No & No & No & Sí \\
Configuración mínima & Sí & Sí & Parcial & No \\
\hline
\end{tabular}
\caption{Comparativa de proveedores de modelos de lenguaje.}
\label{tab:alt-proveedor}
\end{figure}

\textbf{Justificación de la elección.} GitHub Models combina coste cero, compatibilidad directa con la API de OpenAI (sin cambios de código respecto al acceso directo) y calidad de seguimiento de instrucciones equivalente a los modelos de OpenAI, dado que proporciona acceso a los mismos modelos subyacentes. Para un prototipo académico con iteraciones frecuentes, esta combinación resulta óptima. La arquitectura del sistema encapsula el proveedor en una factoría de modelos (\texttt{utils/llm.py}), de modo que la migración a OpenAI directo, Anthropic u otro proveedor compatible con la especificación de API de OpenAI requiere únicamente modificar las variables de entorno, sin cambios en el código de los agentes.


\subsection{Alternativas al almacén de experiencias}
\label{sec:alt-experiencia}

El \emph{pool} de experiencias requiere almacenar registros estructurados con metadatos de ejecución, resultados por modelo candidato y referencias a artefactos MLflow, y recuperarlos eficientemente por similitud de perfil de dataset. La elección del sistema de almacenamiento influye en la complejidad de despliegue, la capacidad de consulta y la integración con el resto del stack.

\textbf{Bases de datos vectoriales} (Chroma, Weaviate, Pinecone) son la opción natural cuando la similitud se calcula sobre embeddings de texto. Sin embargo, la similitud en este sistema se calcula sobre el \texttt{DatasetProfile}, un conjunto de características discretas y categóricas (tamaño del dataset, número de columnas, tipo de problema, presencia de valores ausentes) que no se beneficia de la codificación vectorial continua. Un sistema de recuperación basado en coincidencia de campos categóricos sobre SQL es más interpretable, más rápido para este caso y elimina la dependencia de un servicio externo o de un modelo de embeddings.

\textbf{PostgreSQL / MySQL} ofrecen capacidades relacionales completas con índices, transacciones y soporte para consultas complejas. Son la opción correcta en un sistema de producción con múltiples usuarios concurrentes y millones de registros. Para un prototipo con centenares de registros y un único proceso de acceso, añaden la complejidad de desplegar y configurar un servidor de base de datos, lo que contradice el objetivo de reproducibilidad con \texttt{docker compose up}.

\textbf{MLflow como almacén de metadatos de experiencia} podría parecer una opción natural dado que MLflow ya registra las métricas y los artefactos de cada ejecución. Sin embargo, MLflow no proporciona una API de consulta estructurada sobre atributos personalizados del dataset (como el tipo de problema o el número de filas); su modelo de datos está orientado al seguimiento de parámetros y métricas de experimentos, no a la recuperación por similitud de perfil. Reutilizar MLflow como \emph{pool} de experiencias implicaría sobrecargar su esquema con etiquetas personalizadas y realizar consultas con filtros no indexados, lo que degradaría tanto el rendimiento como la mantenibilidad.

\textbf{Ficheros JSON sin base de datos} son adecuados como copia de auditoría legible por humanos (uso que el sistema ya hace con \texttt{experience\_pool/<task\_id>.json}), pero no ofrecen capacidad de consulta indexada. Recuperar las $k$ experiencias más similares requeriría cargar todos los ficheros en memoria y ordenarlos, lo que no escala incluso para cientos de registros.

\begin{figure}[H]
\centering
\begin{tabular}{lcccc}
\hline
\textbf{Criterio} & \textbf{SQLite} & \textbf{PostgreSQL} & \textbf{BD vectorial} & \textbf{JSON} \\
\hline
Sin servidor externo & Sí & No & Parcial & Sí \\
Consultas relacionales estructuradas & Sí & Sí & No & No \\
Recuperación por perfil categórico & Sí & Sí & Parcial & No \\
Cero configuración de despliegue & Sí & No & No & Sí \\
Esquema jerárquico (3 tablas) & Sí & Sí & No & Parcial \\
\hline
\end{tabular}
\caption{Comparativa de alternativas para el almacén de experiencias.}
\label{tab:alt-experiencia}
\end{figure}

\textbf{Justificación de la elección.} SQLite satisface todos los requisitos del prototipo: es embebida (sin servidor), soporta el esquema relacional de tres tablas, permite consultas SQL indexadas sobre los campos del perfil y se despliega automáticamente junto con el resto del sistema sin configuración adicional. Su limitación de acceso concurrente es irrelevante para un pipeline que ejecuta un único experimento por proceso.


\subsection{Alternativas al sistema de seguimiento de experimentos}
\label{sec:alt-mlflow}

El seguimiento de experimentos cubre el registro de parámetros, métricas y artefactos durante el entrenamiento, y el registro y versionado de modelos en un repositorio central. Esta funcionalidad es indispensable para la reproducibilidad del pipeline y para la comparación de candidatos en el módulo de evaluación.

\textbf{Weights \& Biases (W\&B)} es una plataforma de seguimiento de experimentos ampliamente adoptada en la comunidad de investigación, con una interfaz web de alta calidad y funcionalidades avanzadas de visualización y comparación de modelos~\cite{wandb}. Su limitación principal para este proyecto es que el nivel gratuito es cloud-first: los datos de experimentos se almacenan en los servidores de W\&B, lo que introduce una dependencia de internet durante la ejecución del pipeline y suscita interrogantes de privacidad cuando los datos de entrenamiento son sensibles. Una instalación auto-hospedada (\emph{W\&B Server}) está disponible, pero añade complejidad de despliegue significativa.

\textbf{Neptune.ai} y \textbf{Comet ML} presentan características similares a W\&B: excelente interfaz, API Python madura, pero dependencia cloud en los niveles accesibles para un proyecto académico.

\textbf{DVC (Data Version Control)} gestiona el versionado de datasets y modelos mediante Git, con almacenamiento en S3, GCS u otros backends. Es la herramienta adecuada cuando la reproducibilidad a nivel de datos es el requisito principal. Sin embargo, DVC no está orientado al registro en tiempo real de métricas durante el entrenamiento ni dispone de un Model Registry con etapas de ciclo de vida (Staging, Production, Archived), funcionalidad que el módulo de evaluación del sistema utiliza para recuperar el campeón actual.

\textbf{ClearML} ofrece una plataforma MLOps más completa que MLflow, con orquestación de experimentos, gestión de datasets y pipelines automatizados. Para este proyecto su alcance es excesivo: se necesita el seguimiento de experimentos y un registro de modelos, no una plataforma de orquestación completa que competiría con LangGraph.

\begin{figure}[H]
\centering
\begin{tabular}{lcccc}
\hline
\textbf{Criterio} & \textbf{MLflow} & \textbf{W\&B} & \textbf{DVC} & \textbf{ClearML} \\
\hline
Autohospedado sin configuración compleja & Sí & Parcial & Sí & Parcial \\
Model Registry con etapas de ciclo de vida & Sí & Sí & No & Sí \\
Imagen Docker oficial & Sí & No & No & Sí \\
API Python para logging en tiempo real & Sí & Sí & No & Sí \\
Integración nativa con scikit-learn / LightGBM & Sí & Sí & No & Sí \\
Sin dependencia de servicio cloud externo & Sí & Parcial & Sí & Parcial \\
\hline
\end{tabular}
\caption{Comparativa de herramientas de seguimiento de experimentos.}
\label{tab:alt-mlflow}
\end{figure}

\textbf{Justificación de la elección.} MLflow combina imagen Docker oficial, despliegue auto-hospedado sin dependencia de servicios externos, Model Registry con etapas de ciclo de vida (necesario para que el módulo de evaluación recupere el campeón actual), y la integración más madura con las bibliotecas de modelado utilizadas (scikit-learn, LightGBM, XGBoost, CatBoost). La interfaz web integrada permite al operador inspeccionar los experimentos durante el desarrollo sin necesidad de configurar una herramienta adicional.


\subsection{Alternativas al stack de interfaz de usuario}
\label{sec:alt-frontend}

La interfaz de usuario tiene requisitos específicos que van más allá de una visualización estática de resultados: debe recibir y mostrar eventos en tiempo real a medida que el pipeline se ejecuta (razonamiento del agente, progreso del entrenamiento, mensajes del planificador), y debe gestionar los dos flujos de aprobación humana con formularios interactivos y contexto completo.

\textbf{Streamlit} fue la primera implementación del frontend del sistema. Su propuesta de valor es permitir construir aplicaciones de datos con Python puro, sin necesidad de JavaScript ni HTML. Sin embargo, Streamlit presenta limitaciones estructurales para los requisitos de este proyecto: el soporte de \emph{Server-Sent Events} (SSE) para streaming de eventos desde el backend requiere soluciones no idiomáticas; el modelo de re-ejecución completa de la página ante cualquier interacción dificulta mantener el estado de los flujos HITL; y la arquitectura de un único proceso Python hace difícil separar limpiamente el frontend de la lógica del backend.

\textbf{Gradio} ofrece una propuesta aún más simple que Streamlit, orientada principalmente a la demostración de modelos de ML mediante interfaces generadas automáticamente. Sus componentes predefinidos no son suficientes para las vistas de aprobación del pipeline, que requieren mostrar informes estructurados, métricas comparativas y formularios de decisión con comentarios.

\textbf{Dash (Plotly)} es una alternativa Python que introduce componentes reactivos mediante callbacks. Es adecuada para dashboards de visualización de datos, pero su modelo de programación por callbacks para gestionar estados de flujos multi-paso resulta más complejo que el modelo de componentes de React para este caso de uso.

\textbf{Django / Flask con plantillas Jinja} permitirían construir una interfaz server-side rendering sobre el backend Python existente. Sin embargo, el soporte de SSE en estas plataformas requiere configuración adicional (conexiones de larga duración, gestión de threads), y la integración de componentes reactivos para actualizaciones en tiempo real exige incorporar JavaScript de todas formas.

\textbf{Next.js + FastAPI} desacopla completamente el frontend del backend. FastAPI proporciona soporte nativo y asíncrono para SSE, con tipado de endpoints mediante Pydantic. Next.js permite construir componentes React con actualizaciones en tiempo real mediante el consumo nativo del \texttt{EventSource} del navegador. Esta separación también facilita el despliegue independiente de ambos servicios en contenedores distintos.

\begin{figure}[H]
\centering
\begin{tabular}{lccccc}
\hline
\textbf{Criterio} & \textbf{Next.js+FastAPI} & \textbf{Streamlit} & \textbf{Gradio} & \textbf{Dash} \\
\hline
Soporte nativo SSE/streaming & Sí & Parcial & No & Parcial \\
Componentes reactivos interactivos & Sí & Parcial & No & Sí \\
Flujos HITL con estado persistente & Sí & No & No & Parcial \\
Separación limpia frontend/backend & Sí & No & No & No \\
Despliegue en contenedores independientes & Sí & Parcial & Parcial & Parcial \\
\hline
\end{tabular}
\caption{Comparativa de alternativas para la interfaz de usuario.}
\label{tab:alt-frontend}
\end{figure}

\textbf{Justificación de la elección.} Los dos requisitos diferenciadores de la interfaz, el streaming en tiempo real de eventos del pipeline y los flujos interactivos de aprobación humana, requieren una solución con soporte nativo de SSE y componentes reactivos con estado persistente. Next.js + FastAPI es la única alternativa evaluada que satisface ambos requisitos con un modelo de programación natural. La mayor complejidad inicial de configurar un proyecto React frente a Streamlit queda compensada por la eliminación de los compromisos arquitectónicos que Streamlit impondría en las fases más avanzadas del desarrollo.
